% BEGIN GENERATED ARTIFACT: prf:nor-functionally-complete-propositional-logic
\newpage
\phantomsection
\label{prf:nor-functionally-complete-propositional-logic}
\LRAProofFor{thm:nor-functionally-complete-propositional-logic}

\begin{remark*}[Return]
\hyperref[thm:nor-functionally-complete-propositional-logic]{Return to Theorem (NOR is Functionally Complete).}
\end{remark*}

\begin{theorem*}[NOR is Functionally Complete]
The one-connective basis \(\{\downarrow\}\) is functionally complete.
\end{theorem*}

\begin{proof}
\textbf{Professional Standard Proof.}~
TODO: supply the compact proof.
\end{proof}

\begin{proof}
\textbf{Detailed Learning Proof.}~
TODO: supply the detailed learning proof.
\end{proof}

\begin{remark*}[Proof structure]
Show that NOR defines negation and disjunction. Since \(\{\neg,\lor\}\) is
functionally complete, any Boolean function can be represented by first using
that basis and then replacing the defined connectives by NOR-only patterns.
\end{remark*}

\begin{dependencies}
\begin{itemize}
  \item \hyperref[def:nor-connective-propositional-logic]{NOR Connective}
  \item \hyperref[def:nor-only-formula-propositional-logic]{NOR-Only Formula}
  \item \hyperref[def:functionally-complete-connective-set-propositional-logic]{Functionally Complete Connective Set}
  \item \hyperref[thm:negation-disjunction-functionally-complete-propositional-logic]{Negation and Disjunction are Functionally Complete}
\end{itemize}
\end{dependencies}

\clearpage
% END GENERATED ARTIFACT: prf:nor-functionally-complete-propositional-logic
